%% ============================================================
%%  Chapter 2 — Scope and Methodology
%% ============================================================
\chapter{Scope and Methodology}
\label{ch:scope}

\section{Review Scope}

\begin{description}[leftmargin=3cm, style=nextline]
  \item[In scope]
    \fillme{Enumerate what was reviewed -- e.g. source tree, compiled bytecode, config files,
    dependency manifests -- with rough size (file/class counts) if useful context.}
    \begin{itemize}
      \item \fillme{e.g. Views / templates (~N files)}
      \item \fillme{e.g. Web/security XML or annotation-based configuration}
      \item \fillme{e.g. Properties / environment configuration files}
      \item \fillme{e.g. Compiled classes, analyzed via bytecode disassembly (source was absent)}
    \end{itemize}

  \item[Out of scope]
    \fillme{e.g. Dynamic testing; running the application; network calls; binary data files.
    No code was modified.}

  \item[Context path]
    \fillme{\url{https://example.com/app} (used in illustrative PoCs only -- not executed).}
\end{description}

\section{Methodology}

The assessment followed a static, read-only white-box approach:

\begin{enumerate}
  \item \textbf{Configuration analysis} --- \fillme{web/security config, properties files, ORM
    mapping config.}
  \item \textbf{Secret scanning} --- all configuration and selected constant pools scanned for
    credentials, keys, and tokens.
  \item \textbf{Authorization model analysis} --- \fillme{the framework's auth filter chain and
    any public-endpoint allowlist.}
  \item \textbf{Input/output handling analysis} --- \fillme{templates/views scanned for unescaped
    output, raw query construction.}
  \item \textbf{Bytecode / binary analysis} --- \fillme{if source was unavailable, note the
    disassembly tool used and what was searched for (injection sinks, crypto primitives, PRNG use).}
  \item \textbf{Dependency audit} --- library version fingerprinting against known CVE databases.
\end{enumerate}

\section{Confidence Levels}

Each finding is assigned one of two confidence levels:

\begin{description}[leftmargin=3cm]
  \item[\Confirmed] Direct evidence quoted from source or bytecode --- the vulnerability exists
    as described.
  \item[\Suspected] The mechanism is present but the full taint path or runtime-derived data
    could not be verified statically. These require dynamic confirmation (see
    Appendix~\ref{app:runtime}).
\end{description}

\section{Limitations}

\begin{itemize}
  \item \fillme{Note any part of the codebase that could not be verified statically (e.g. runtime
    DB-driven authorization tables, dynamic includes).}
  \item \fillme{Note any deployment/config-dependent finding whose real-world severity depends on
    something only knowable on the live server (see Appendix~\ref{app:runtime}).}
\end{itemize}
